\section{Repository architecture}

The code is organised by role, and \cmd{sim/README.md} is the
authoritative file-by-file map; this section gives the shape.

\begin{description}
  \item[sim/core/] the solver and specimen machinery: the
    label-indexed anisotropic FDTD solver, the pseudospectral
    solver, the specimen builder, rotation machinery, and the
    bit-equivalence gate.
  \item[sim/model/] forward models and inversion: the born ladder,
    the fabric and sweep fits, the gridded inversion and its daemon,
    and the ray model the studio uses.
  \item[sim/pipeline/] drivers: the resumable sweep engine, trace
    lab, arc backprojection and visualisation.
  \item[sim/fe\_crosscheck/] the independent finite-element arbiter
    and its archived cross-check rounds.
  \item[sim/fw\_checks/ and sim/validation/] the validation chain of
    Chapter~\ref{ch:forward}.
  \item[sim/analysis/] roughly 75 analysis modules with their stored
    results, grouped by theme, each self-contained and runnable from
    its own directory.
  \item[sim/ref/, sim/runs/, sim/results/, sim/figures/] reference
    traces, batch campaign drivers, stored campaign results, and the
    paper's figure scripts.
  \item[gui/] the Streamlit studio, the sweep tab, and the launcher
    runtime scripts.
  \item[vendor/] vendored dependencies, so the repository is
    self-contained.
\end{description}

Two conventions hold everywhere: modules read stored inputs relative
to their own location, so they run from their own directory with no
environment setup; and expected outputs are recorded in docstrings,
so checking a module means running it and comparing.
